% concatenated from new4__3_.tex
\documentclass{article}
\usepackage{amsmath, amssymb, amscd, amsthm, amsfonts}
\usepackage{graphicx}
\usepackage{mathrsfs}
\usepackage{hyperref}
\usepackage{color}
\usepackage{float}
\usepackage{epstopdf}
\usepackage{tikz}
\usepackage{caption}
\usepackage{subcaption}
\usepackage{booktabs}
\usetikzlibrary{arrows.meta}
\usetikzlibrary{calc}

\renewcommand{\baselinestretch}{1.2}
\topmargin  = -0.2 in
\oddsidemargin = 0.15 in
\setlength{\textheight}{8.5in}
\setlength{\textwidth}{6in}
\setlength{\unitlength}{1.0 mm}

\newcommand{\Ftwo}{\mathbb{F}_2}
\newcommand{\mchi}{\hat{\chi}}
%\newcommand{\mchi}{\lceil\chi\rceil}
\newcommand{\tw}{{\rm tw}}
\newcommand{\2}{\vspace{0.2cm}}
\newcommand{\fo}{f_{o}}
\newcommand{\comp}[1]{\overline{#1}}

\newtheorem{theorem}{Theorem}[section]
\newtheorem{proposition}[theorem]{Proposition}
\newtheorem{lemma}[theorem]{Lemma}
\newtheorem{remark}[theorem]{Remark}
\newtheorem{corollary}[theorem]{Corollary}
\newtheorem{conjecture}[theorem]{Conjecture}
\newtheorem{claim}[theorem]{Claim}
\newtheorem{question}[theorem]{Question}
\newtheorem{definition}[theorem]{Definition}
\newtheorem{example}[theorem]{Example}


%\title{Generalization of Ferber-Krivelevich Theorem on parity of degrees in induced subgraphs}
\title{Large induced subgraphs with prescribed degree parity}
\author{Jiangdong Ai\thanks{School of Mathematical Sciences and LPMC, Nankai University. {\tt jd@nankai.edu.cn}.}
\hspace{2mm}
Qiwen Guo\thanks{Department of Computer Science, Royal Holloway University of London. {\tt gqwmath@163.com}.}
\hspace{2mm}
Gregory Gutin\thanks{Corresponding author. Department of Computer Science, Royal Holloway University of London,  and School of Mathematical Sciences and LPMC, Nankai University. {\tt g.gutin@rhul.ac.uk}.}
\hspace{2mm} Yiming Hao\thanks{School of Mathematical Sciences and LPMC, Nankai University. {\tt  
1120230031@mail.nankai.edu.cn}.}
\hspace{2mm} Anders Yeo\thanks{Department of Mathematics and Computer Science, University of Southern Denmark, and Department of Mathematics, University of Johannesburg.   {\tt andersyeo@gmail.com}.}
}

\begin{document}

\maketitle

\begin{abstract}
A long-standing conjecture of Caro (Discrete Math, 1994), confirmed by Ferber and Krivelevich (Adv\ Math, 2022), states that every $n$-vertex graph $G$ without isolated vertices contains an induced subgraph of order linear in $n$ in which every vertex has odd degree. We generalize this result to graphs $G$ whose vertices are labeled by $\ell: V(G)\to \{0,1\}$. We require, in an induced subgraph, all
$0$-labeled vertices to have even degree and all $1$-labeled vertices to have odd degree. Let $h_{\ell}(G)$ denote the maximum order of such a subgraph. 
Let $f_{oe}(G)=\min_{\ell} h_{\ell}(G)$ be the worst-labeling parameter. We establish a pointwise lower bound for $h_{\ell}(G)$ that immediately yields a linear lower bound in $|V(G)|$ for $f_{oe}(G)$, where $G$ has no isolated vertices. 
For an $n$-vertex connected graph, we obtain a sharp lower bound for $f_{oe}(G)$: 
$f_{oe}(G)\ge \lceil (n-1)/\mchi(G) \rceil ,$ where $\mchi(G)$ is the maximum chromatic number of a minor of $G.$ Using proved cases of Hadwiger's Conjecture, we show that for $t\in \{3,4,5,6\}$, if an $n$-vertex connected graph $G$ is $K_t$-minor-free, then $f_{oe}(G)\ge \lceil (n-1)/(t-1)\rceil$ and this bound is sharp for each $t\in \{3,4,5,6\}$. 
Finally, we conjecture that $f_{oe}(G)\ge f_o(G)/2$ for all graphs $G$ and confirm the conjecture for all trees and complete multipartite graphs.
\end{abstract}

\section{Introduction}

%For a graph $G$, an \emph{odd induced subgraph} is an induced subgraph in which every vertex has odd degree. 
A graph $H$ is {\em odd} if the degree of every vertex of $H$ is odd. A key result on odd spanning subgraphs was obtained by 
Scott \cite{Scott2001} who proved that every connected graph $G$ with even number of vertices has an odd spanning forest $F$ such that
each component of $F$ is an induced subgraph of $G$. %Clearly, such spanning forests generalize perfect matchings and
%thus Gutin \cite{Gutin2016} coined them {\em perfect forests}. 
Gutin \cite{Gutin2016} and  Caro, Lauri and Zarb \cite{CLZ} found short proofs of
Scott's theorem. Gutin and Yeo \cite{GutinYeo2022} generalized Scott's theorem to graphs whose vertices are labeled by parities such that the sum of the parities is even.
We call a labeling $\ell:\ V(G)\to \{0,1\}$ {\em even-sum} if $\sum_{v\in V(G)}\ell(v)$ is even. 


\begin{theorem}\label{thm:GY}
If a graph $G$ is connected and a labeling $\ell:\ V(G)\to \{0,1\}$ is even-sum, 
then $G$ has a spanning forest $F$ such that every component of $F$ is an induced subgraph of $G$ and 
the degree $d_F(x)\equiv \ell(x) {\pmod 2}$ for every $x\in V(F).$ 
\end{theorem}

Henceforth, an induced subgraph $H$ of $G$ is called {\em $\ell$-admissible} if $d_H(x)\equiv \ell(x) {\pmod 2}$ for every $x\in V(H).$
We also call $V(H)$ an {\em $\ell$-admissible} set. 

For odd induced subgraphs, let $f_o(G)$ denote the maximum order of an odd induced subgraph of a graph $G$.
Let $G$ be a graph of order $n$ without isolated vertices. Resolving a problem of Noga Alon, Caro \cite{Caro1994} proved that $f_o(G)=\Omega(\sqrt{n})$ and conjectured the linear lower bound $f_o(G)=\Omega(n)$. Scott \cite{Scott1992} subsequently obtained $f_o(G)=\Omega(n/\log n)$, and the conjecture was finally settled by Ferber and Krivelevich.

\begin{theorem}[Ferber-Krivelevich \cite{FK}]\label{thm: os linear bound}
Every graph of order $n$ without isolated vertices contains an odd induced subgraph of order at least $cn$, where $c=10^{-4}$.
\end{theorem}

In Section~\ref{sec:FK} we generalize this theorem to graphs whose vertices are labeled by parities. 
Given a labeling $\ell:\ V(G)\to\{0,1\}$, write $V_i=\ell^{-1}(i)$ and define
$$
h_{\ell}(G):=\max\bigl\{\,|U|:\ U\subseteq V(G), U \mbox{ is $\ell$-admissible}\bigr\}.
$$
We also introduce the worst-labeling parameter
$$
f_{oe}(G):=\min_{\ell: V(G)\to\{0,1\}} h_{\ell}(G),
$$
which measures the largest $\ell$-admissible induced subgraph guaranteed against an adversarial labeling. 
%Thus an admissible induced subgraph $G[U]$ has even degree at every $0$-labeled vertex of $U$ and odd degree at every $1$-labeled vertex of $U$. 
Our main result, Theorem~\ref{oddevenIS}, is a point-wise lower bound on $h_{\ell}(G)$ expressed in terms of $|V_0|$, the number of isolated vertices of $G[V_1]$, and $f_o(G[V_1])$. Combined with Theorem~\ref{thm: os linear bound}, it yields the universal linear bound $f_{oe}(G)\ge |V(G)|/10004$ (Corollary~\ref{cor:universal}).


To prove Theorem~\ref{oddevenIS}, we will also use the following
classical theorem of Gallai (\cite{L1993}, Problem 5.17). 

\begin{theorem}[Gallai's Theorem]\label{thm: gallai thm}
For every graph $G$, the vertex set $V(G)$ admits
\begin{enumerate}
    \item a partition $V(G)=V_1\cup V_2$ such that all degrees in both $G[V_1]$ and $G[V_2]$ are even;
    \item a partition $V(G)=V_1\cup V_2$ such that all degrees in $G[V_1]$ are odd and all degrees in $G[V_2]$ are even.
\end{enumerate}
\end{theorem}

\medskip


Scott \cite{Scott1992} proved that $f_o(G)\ge |V(G)|/(2\chi(G))$ for every graph $G$ without isolated vertices. He posed the following:
\begin{conjecture}\label{conj:scott}
For every graph $G$ without isolated vertices, $f_o(G)\ge |V(G)|/\chi(G).$   
\end{conjecture}
Wang and Wu \cite{WW} showed that Conjecture \ref{conj:scott} fails for bipartite graphs, but holds for line graphs. It is still unknown whether the conjecture holds for 
non-bipartite graphs.
Very recently, Ning \cite{Ning} extended the positive result in \cite{WW} by confirming Conjecture \ref{conj:scott} for claw-free graphs without isolated vertices. 

In this paper, we will use the following parameter related to $\chi(G)$ and introduced by Halin \cite{Halin} who called it {\em modified chromatic number}.
For a graph $G$, let
$$
  \mchi(G)=\max\{\chi(H): H \text{ is a minor of } G\}
$$
Barioli et al. \cite{Barioli+} extended this definition to other graph parameters. 
%and called them {\em minor-monotone-ceiling parameters}. 
Hogben \cite{Hogben} is an informative survey on the topic. 


%Unfortunately, $\mchi(G)$ is not bounded by any function of $\chi(G).$ For example, $\chi(K_{t,t})=2$ for every $t,$ but $\mchi(K_{t,t})\ge t+1$ as $K_{t,t}$ contains $K_{t+1}$ as a minor (contract $t-1$ disjoint edges).  
The main result of Section \ref{sec:trees} is Theorem \ref{thm:tw} whose proof uses Theorem \ref{thm:GY} thus providing a novel link between spanning and induced $\ell$-admissible subgraphs. The main part of Theorem \ref{thm:tw}
is the following sharp bound on $f_{oe}(G)$ for every connected graph $G$ of order $n$: $f_{oe}(G)\ge \lceil (n-1)/\mchi(G) \rceil$. 
This bound can be useful for classes $\cal C$ of connected graphs $G$, where $\mchi(G)$ is bounded by a constant. We show that $\mchi$ is bounded in $\cal C$ by a constant if and only if $\cal C$ excludes some fixed complete graph $K_t$ as a minor. 

Using proved cases of Hadwiger's Conjecture, we show that for $t\in \{3,4,5,6\}$, if an $n$-vertex connected graph $G$ is $K_t$-minor-free, then $f_{oe}(G)\ge \lceil (n-1)/(t-1)\rceil$ and this bound is sharp for every
$t\in \{3,4,5,6\}$. In particular, for connected $K_4$-minor-free graphs, we have  $f_{oe}(G)\ge \lceil (n-1)/3\rceil .$ It is different from the bound
$f_o(G)\ge 2n/5$ for $K_4$-minor-free graphs without isolated vertices proved by Hou et al. \cite{HYLL}.
Also, for a connected planar $n$-vertex graph $G$, we have $f_{oe}(G)\ge \lceil (n-1)/4 \rceil$ and this bound is sharp. Indeed, planar graphs are $K_5$-minor-free and the sharpness is due to $K_{2,2,2}$, see the case of $t=5$ in Theorem \ref{thm:Hadw}. As far as we know, there was no sharp lower bound for $f_o(G)$, where $G$ is a connected planar graph. The bound $f_o(G)\ge n/(2\chi(G))$ of \cite{Scott1992} implies only that $f_o(G)\ge n/8$. 
%However, if a planar graph $G$ without isolated vertices has girth $g(G)\ge 6$, then $f_o(G)\ge n/3$ \cite{RHZ}. Ning \cite{Ning} suspects that if $G$ is also bipartite, then $f_o(G)\ge n/2.$

 
The second part of Theorem \ref{thm:tw} is the claim that for every connected graph $G$ of order $n$ and an even-sum labeling $\ell$,
we have $h_{\ell}(G)\ge \lceil n/\mchi(G) \rceil$. This, in particular, allows us to show that every connected chordal graph with even number of vertices satisfies Conjecture \ref{conj:scott}.
It would be interesting to know whether every connected chordal graph with odd number of vertices satisfies Conjecture \ref{conj:scott}.


In Section \ref{sec:conj}, we conjecture that $f_{oe}(G)\ge f_o(G)/2$ for all graphs $G$ and prove the inequality for all trees and complete multipartite graphs.


%Gallai's proof in \cite{L1993} is purely combinatorial. The natural digraph analogue, where ``degree'' is replaced by ``out-degree'', does not always hold: call a partition $V=V_0\cup V_1$ of a digraph $D=(V,A)$ \emph{even-even} if every vertex of $D[V_0]$ and $D[V_1]$ has even out-degree, and \emph{even-odd} if every vertex of $D[V_0]$ has even out-degree while every vertex of $D[V_1]$ has odd out-degree. The directed $3$-cycle admits no even-even partition, and the directed $2$-vertex path admits no even-odd partition.

%In Section~\ref{sec:Gallai} we give a complete characterization of digraphs admitting such partitions via a linear algebraic approach over $\Ftwo$, in the spirit of \cite{Gutin2016}. We unify the even-even and even-odd cases by allowing each vertex to prescribe a possibly different parity for each side; the resulting Theorem~\ref{thm:general-digraph} reduces the existence question to the solvability of a single linear system over $\Ftwo$ and contains both classical cases as immediate corollaries. The same system also yields an exact count of all valid partitions (Corollary~\ref{cor:counting}), structural results for bipartite and symmetric digraphs (Proposition~\ref{prop:bipartite-symmetric}), and a complete combinatorial characterization for functional digraphs (Proposition~\ref{prop:functional}). All algorithms run in $O(|V|^3)$ arithmetic operations over $\Ftwo$ via Gaussian elimination.

\section{Generalizing Ferber--Krivelevich Theorem}\label{sec:FK}

We use the following theorem of Scott \cite{Scott1992}.

\begin{theorem}[Scott \cite{Scott1992}]\label{thm:Scott}
Let $G$ be a graph without isolated vertices and with a maximum independent set of size $p$. Then $G$ contains an odd induced subgraph of order at least $p/2$.
\end{theorem}

We extend the notation $f_o(H)$ to arbitrary graphs $H$ (possibly with isolated vertices) by defining it as the maximum order of an odd induced subgraph of $H$; isolated vertices are excluded since no odd induced subgraph can contain one.

The main result of this section is the following pointwise lower bound.

\begin{theorem}\label{oddevenIS}
Let $G=(V,E)$ be a graph without isolated vertices and let $\ell:\ V\to\{0,1\}$. Set $V_i=\ell^{-1}(i)$ for $i\in\{0,1\}$, and let $I$ be the set of isolated vertices of $G[V_1]$. Then
$$
h_{\ell}(G)\ge \max\Bigl\{\frac{|V_0|}{2},\ \frac{|I|}{2},\ f_o(G[V_1])\Bigr\}.
$$
\end{theorem}

\begin{proof}
We give three constructions.

\medskip
\noindent\textbf{Construction A: use the $0$-labeled vertices.}
Apply Theorem~\ref{thm: gallai thm}(1) to the induced subgraph $G[V_0]$. We obtain a partition $V_0=A\cup B$ such that both $G[A]$ and $G[B]$ have all degrees even. One part has size at least $|V_0|/2$; assume $|B|\ge|A|$. Then every vertex of $G[B]$ has even degree and every vertex of $B$ is labeled 0, so
$$
h_{\ell}(G)\ge|B|\ge\frac{|V_0|}{2}.
$$

\medskip
\noindent\textbf{Construction B: use the $1$-labeled vertices.}
Any odd induced subgraph of $G[V_1]$ is $\ell$-admissible in $G$ and so $h_{\ell}(G)\ge f_o(G[V_1])$.

\medskip
\noindent\textbf{Construction C: use the isolated vertices of $G[V_1]$.}
If $I=\varnothing$ the bound $|I|/2=0$ is trivial. So assume $I\neq\varnothing$. Consider $G'':=G[V_0\cup I]$. Every $u\in I$ has a neighbor in $V_0$ (since $G$ has no isolated vertices and vertices of $I$ have no neighbors in $V_1$); also $I$ is independent in $G''$.

Let $D\subseteq V_0$ be an inclusion-minimal subset that dominates $I$ in $G''$. We have the following: %By minimality, for each $w\in D$ there is a ``private'' neighbor $u_w\in I$ with $N(u_w)\cap D=\{w\}$.

\begin{claim}\label{clm:private}
For every $w\in D$ there exists $u_w\in I$ with $N_{G''}(u_w)\cap D=\{w\}$.
\end{claim}
\begin{proof}
If no such $u_w$ existed, every neighbor of $w$ in $I$ would have another neighbor in $D\setminus\{w\}$, so $D\setminus\{w\}$ would still dominate $I$, contradicting minimality.
\end{proof}

Let $I_D:=\{u_w:w\in D\}$. Choose $D'\subseteq D$ uniformly at random. Define
$$
I_0:=\{u\in I\setminus I_D:\ |N(u)\cap D'|\text{ is odd}\}.
$$
For each $u\in I\setminus I_D$, fix $w_u\in N(u)\cap D$; the involution $S\mapsto S\triangle\{w_u\}$ shows that exactly half of all subsets of $D$ give odd intersection with $N(u)$, so
$$
\mathbb{P}\bigl(|N(u)\cap D'|\text{ is odd}\bigr)=\tfrac{1}{2}.
$$
Define
$$
I_1:=\{u_w\in I_D:\ w\in D'\text{ and }\deg_{G[D'\cup I_0]}(w)\text{ is odd}\},
$$
and set $U:=D'\cup I_0\cup I_1$.

\emph{Parity verification.} For each $u\in I_0$, all neighbors of $u$ that lie in $U$ belong to $D'$ (since $I$ is independent), 
so $\deg_{G[U]}(u)=|N(u)\cap D'|$, which is odd by definition. Each $u_w\in I_1$ has $N(u_w)\cap D\cap U=\{w\}$ and 
no neighbors in $I$, so $\deg_{G[U]}(u_w)=1$, which is odd. For $w\in D'\subseteq V_0$, if $u_x\in I_1$ 
and $x\neq w$ then $w\notin N(u_x)\cap D$ by Claim~\ref{clm:private}, so
$$
\deg_{G[U]}(w)\equiv \deg_{G[D'\cup I_0]}(w)+\mathbf{1}_{\{u_w\in I_1\}}\pmod{2},
$$
and the definition of $I_1$ ensures that this is even. Thus, $\deg_{G[U]}(v)\equiv \ell(v)\pmod{2}$ for all $v\in U$.

\emph{Size bound.} By linearity of expectation,
$$
\mathbb{E}|U|\ge\mathbb{E}|D'|+\mathbb{E}|I_0|=\frac{|D|}{2}+\frac{|I\setminus I_D|}{2}\ge
\frac{|D|}{2}+\frac{|I|-|D|}{2}=\frac{|I|}{2}.
$$
Hence, there exists a choice of $U$ with $|U|\ge |I|/2$, giving $h_{\ell}(G)\ge|I|/2$.

Combining Constructions A, B, and C completes the proof.
\end{proof}

The point-wise theorem yields the following corollaries.

\begin{corollary}\label{cor:Scott-bound}
Under the assumptions of Theorem~\ref{oddevenIS},
$$
h_{\ell}(G)\ge \frac{1}{2}\max\Bigl\{|V_0|,\ |I|,\ \alpha(G[V_1]-I)\Bigr\},
$$
where $\alpha(\cdot)$ denotes the independence number.
\end{corollary}
\begin{proof}
By Theorem~\ref{thm:Scott}, $G[V_1]-I$ has an odd induced subgraph of order at least $\alpha(G[V_1]-I)/2$. Since $f_o(G[V_1])=f_o(G[V_1]-I)$, the result follows from Theorem~\ref{oddevenIS}.
\end{proof}

\begin{corollary}\label{cor:universal}
Suppose every graph $H$ without isolated vertices satisfies $f_o(H)\ge c|V(H)|$ for some $c>0$. Then every graph $G$ without isolated vertices and every labeling $\ell:\ V(G)\to\{0,1\}$ satisfy
$$
h_{\ell}(G)\ge \frac{c}{1+4c}\,|V(G)|.
$$
In particular, Theorem~\ref{thm: os linear bound} implies $f_{oe}(G)\ge \frac{1}{10004}\,|V(G)|$.
\end{corollary}
\begin{proof}
Let $n=|V(G)|$, $x=|V_0|/n$, $y=|I|/n$, and $\kappa=c/(1+4c)$. Since $G[V_1]-I$ has no isolated vertices, $f_o(G[V_1])\ge c(1-x-y)n$. Suppose for contradiction that all three quantities in Theorem~\ref{oddevenIS} are less than $\kappa n$. Then $x<2\kappa$, $y<2\kappa$, and $1-x-y<\kappa/c$. Adding the three inequalities gives us $1<4\kappa+\kappa/c=(4c+1)/(1+4c)=1$, a contradiction.
\end{proof}

\section{Bounding $h_{\ell}(G)$ for connected graphs using modified chromatic number}\label{sec:trees}

Recall that $\mchi(G)=\max\{\chi(H): H \text{ is a minor of } G\}.$

%Let $\tw(G)$ denote the tree-width of a graph $G.$ The main result of this section is follows.

\begin{theorem}\label{thm:tw}
Let $G$ be a connected graph of order $n$.
%Then $V(G)$ can be partitioned into at most $\tw(G)+1$ sets such that each subgraph of $G$ induced by such a set is  $\ell$-admissible. 
Then \begin{equation}\label{ineq:1} f_{oe}(G)\ge  \lceil (n-1)/\mchi(G) \rceil\end{equation}  
If a labeling $\ell$ of $V(G)$ is even-sum, then \begin{equation}\label{ineq:2}h_{\ell}(G)\ge  \lceil n/\mchi(G) \rceil\end{equation}
Both bounds are sharp.
\end{theorem} 
\begin{proof}
Let $\ell:\ V(G)\to\{0,1\}$ be an arbitrary labeling.  Consider two cases.

\2

\noindent {\bf Case 1: $\ell$ is even-sum.} By Theorem \ref{thm:GY}, $V(G)$ can be partitioned into sets $X_1,\dots ,X_p$ such that each $G[X_i]$ is connected and $\ell$-admissible.  We will call the subgraphs $G[X_i]$ {\em $\ell$-subgraphs}. Let $G^*$ be obtained from $G$ by contracting $X_1,\dots ,X_p$ and let $x_1,\dots ,x_p$ be the corresponding vertices of $G^*.$ Let $c$ be a proper coloring of $G^*$ into $\chi(G^*)$ colors and 
let the color of $\ell$-subgraph $G[X_i]$ be the color of $x_i$ in $c$ for each $i\in [p].$ Observe that there is no edge between any $\ell$-subgraphs of the same color. 
Thus, the union of all $\ell$-subgraphs of the same color is $\ell$-admissible. 

Hence, $V(G)$ can be partitioned into $\chi(G^*)$ sets each of which induces an $\ell$-admissible subgraph. Let $s$ be the maximum size of such a set. Clearly, $s\ge \lceil n/\chi(G^*) \rceil\ge \lceil n/\mchi(G)\rceil.$ 
%Since tree-width is minor-monotone \cite{Dies}, we have $\tw(G^*)\le \tw(G).$  Since for every graph $H$,  $\chi(H)\le \tw(H)+1$ \cite{Dies}, we have $\chi(G^*)\le \tw(G^*)+1\le \tw(G)+1.$ 
Thus, (\ref{ineq:2}) holds.
%\begin{equation}\label{ineq1} h_{\ell}(G)\ge \lceil n/\mchi(G)\rceil\end{equation}

\2

\noindent {\bf Case 2: $\ell$ is not even-sum.} Let $U:=\ell^{-1}(1)=\{v\in V(G): \ell(v)=1\}$ and let $T$ be a spanning tree of $G$. We will show that there is a vertex $u\in U$ such that every connected component of $T-u$ contains an even number of vertices from $U$.
If $|U|=1$, take $u$ to be the unique vertex in $U$. Then every connected component of $T-u$ contains no vertex from $U$. Now suppose that $|U|\ge 3$. Let $T_U$ be the smallest connected subtree of $T$ containing $U$. Choose any leaf $u$ of $T_U.$
By minimality of $T_U$, every leaf of $T_U$ must belong to $U$. Hence $u\in U$.
Since $u$ is a leaf of $T_U$, among the connected components of $T-u$, at most one component contains vertices from $U\setminus\{u\}$. In fact, that component contains all of $U\setminus\{u\}$, while the other components contain no vertices from $U$. Thus, the numbers of vertices from $U$ in the connected components of $T-u$ are
\(
|U|-1,\ 0,\ 0,\ldots,0.
\)
Therefore, every connected component $C$ of $T-u$ satisfies
\(
\sum_{v\in V(C)} \ell(v)\equiv 0\pmod 2.
\)

Let $C'_1,\dots ,C'_q$ be connected components of $G-u$. 
Observe that for each $i\in [q]$, $V(C'_i)$ is the union of the vertex sets of some connected components of $T-u$.
Thus, we have \(
\sum_{v\in V(C_i')} \ell(v)\equiv 0\pmod 2
\) for each $i\in [q]$. 
We can apply Theorem \ref{thm:GY} to each connected component of $G-u$ and conclude that for each $i\in [q]$, as in Case 1,
there is a set $U_i\subseteq V(C'_i)$ such that $G[U_i]$ is $\ell$-admissible in $C'_i$ and $|U_i|\ge \lceil |V(C'_i)|/\mchi(C'_i) \rceil.$
Observe that $G[U_1\cup \dots \cup U_q]$ is $\ell$-admissible in $G$.
Since $\mchi(C'_i)\le \mchi(G)$, we have $$|U_1\cup \dots \cup U_q|\ge \sum_{i=1}^q \lceil |V(C'_i)|/\mchi(G) \rceil\ge \lceil (n-1)/\mchi(G)\rceil.$$ Thus, $h_{\ell}(G)\ge \lceil (n-1)/\mchi(G)\rceil .$

Since $\ell$ is arbitrary, (\ref{ineq:1}) follows from the above bound and (\ref{ineq:2}). 

To see that both bounds are sharp consider $K_2$ with both vertices labeled by 0. 
\end{proof}

For each of the two bounds of Theorem \ref{thm:tw} there is a number of interesting consequences. We will start with bound (\ref{ineq:1}) and show how to use 
it to obtain a bound of the form $f_{oe}(G)\ge (n-1)/c,$ where $c$ is a constant. We can use the following:

\begin{proposition}\label{prop:finite}
Let $\cal C$ be a class of graphs. Then $\mchi$ is bounded in $\cal C$ by a constant if and only if $\cal C$ excludes some fixed complete graph $K_t$ as a minor.    
\end{proposition}
\begin{proof}
If $\mchi(G)\le k$ then $G$ cannot have a $K_{k+1}$ minor as $\chi(K_{k+1})=k+1.$ Conversely, if every $G\in \cal C$ is $K_t$-minor-free, 
then every minor of every $G\in \cal C$ is also $K_t$-minor-free. Delcourt and Postle \cite{DP} proved that every $K_t$-minor-free graph has chromatic number bounded by 
$O(t\log\log t).$ Hence, $\mchi(G)=O(t\log\log t).$
\end{proof}


Unfortunately, Proposition \ref{prop:finite} does not give us the optimal value of $c$ in $f_{oe}(G)\ge (n-1)/c.$ To get it, we can use
Hadwiger's Conjecture: $\chi(G)\le \eta(G)$ for each graph $G$, where $\eta(G)$ is the maximum $t$ such that $K_t$ is a minor of $G$. Observe that Hadwiger's Conjecture is equivalent to $\mchi(G)=\eta(G)$ for every $G$. Indeed, $K_{\eta(G)}$ is a minor of $G$ so $\mchi(G)\ge \eta(G).$ Conversely, if Hadwiger's Conjecture holds, then every minor $H$ of $G$ satisfies $\chi(H)\le \eta(H)\le \eta(G).$ Thus, $\mchi(G)\le \eta(G).$ Hadwiger's Conjecture was proved for $\eta(G)\le 5$ \cite{Dies}. Thus, we have the following:

\begin{theorem}\label{thm:Hadw}
For $t\in \{3,4,5,6\}$, if an $n$-vertex connected graph $G$ is $K_t$-minor-free, then $f_{oe}(G)\ge \lceil (n-1)/(t-1)\rceil$ and this bound is sharp. 
\end{theorem}
\begin{proof}
Since the bound itself was proved before this theorem, it is enough to prove sharpness.

\2

{\bf $t=3$ (trees):} Observe that the bound becomes an equality already for $K_2$; label both vertices by 0. 

\2 

{\bf $t=4$ (partial 2-trees):} Observe that the bound becomes an equality already for $K_3$ with labels 0, 0, 1. 

\2 

{\bf $t=5$ ($K_5$-minor-free graphs):} The bound becomes an equality for $K_{2,2,2}$. Indeed, $f_{oe}(K_{2,2,2})\ge 2$ and $f_o(K_{2,2,2})=2.$ To show that $f_o(K_{2,2,2})=2$, consider a maximum-order induced odd subgraph $H$ of $K_{2,2,2}$ and let $x_1,x_2,x_3$ be the number of vertices of $H$ in each partite set of $K_{2,2,2}$. If $\min\{x_1,x_2,x_3\}=0$ then 
we observe that $f_o(K_{2,2})=2.$ Otherwise, for each $i\in [3]$, we have $x_i>0$ and $|V(H)|-x_i$ is odd. Now, if $|V(H)|$ is odd,
then each $x_i=2$, a contradiction with $|V(H)|$ being odd. If $|V(H)|$ is even then each $x_i=1$, a contradiction with $|V(H)|$ being even.
Note that $K_{2,2,2}$ is $K_5$-minor-free as it is planar. 

\2

{\bf $t=6$ ($K_6$-minor-free graphs):} Let $C_7=v_0v_1v_2\dots v_6v_0$ and consider $G=\comp{C_7}.$ We will prove that $G$ is $K_6$-minor-free and $f_o(G)=2=\lceil 6/5\rceil.$
Suppose that $G$ has a $K_6$ minor. Since $G$ has only
seven vertices, a $K_6$-minor model either has six singleton branch sets,
which would give a proper $K_6$ subgraph, or has one branch set of size $2$ and five
singleton branch sets.  In the latter case, the five singleton branch sets
would have to be pairwise adjacent, giving a $K_5$ subgraph.  This is
impossible because the maximum clique $\omega(G)=3$.  Therefore $G$ is $K_6$-minor-free.

Now we prove that $f_o(G)=2$. Let $G[S]$ be an odd subgraph. Clearly, $f_o(G)\ge 2$ and 
since $|S|$ is even, to complete the proof we should exclude the cases $|S|=4$ and $|S|=6$. 
For each $v\in S$,
\[
        d_{G[S]}(v)=|S|-1-d_{C_7[S]}(v).
\]
If $|S|=4$ or $|S|=6$, then $|S|-1$ is odd.  Since $d_{G[S]}(v)$ must
be odd, it follows that $d_{C_7[S]}(v)$ is even for every $v\in S$.
But $C_7[S]$ is a proper induced subgraph of a cycle, hence is a forest.  A
forest in which every vertex has even degree has no edges.  Consequently $S$
would have to be an independent set in $C_7$.  This is impossible for
$|S|=4$ or $|S|=6$, because the largest independent sets of $C_7$ have only 3 vertices each.
\end{proof}



Bound (\ref{ineq:2}) implies $f_o(G)\ge  \lceil n/\mchi(G) \rceil$ for a connected graph $G$ with even number of vertices. We will show that this bound can be used for constructing graph classes satisfying Conjecture \ref{conj:scott}. Recall that Conjecture \ref{conj:scott} claims that $f_o(G)\ge n/\chi(G)$ for every graph $G$ without isolated vertices and it was proved for claw-free graphs without isolated vertices \cite{Ning}. Hence, we are interested in graph classes that include some non-claw-free graphs. %The main idea is to consider graph classes for which $\mchi=\chi$.

\begin{corollary}\label{cor:Scott}
Let $\cal C$ be a class of connected graphs with even number of vertices
such that for every graph $G$ in $\cal C$ we have $\mchi(G)=\chi(G)$. 
Then all graphs in $\cal C$ satisfy Conjecture \ref{conj:scott}.    
\end{corollary}
%\begin{proof}
%Let an $n$-vertex graph $G\in \cal C$ and let $N$ be a component of $G$. Then  $f_o(N)\ge n_N/\chi(N)$, where $n_N$ is the order of $N.$ Let ${\cal P}$ be the set of components of $G$. 
%We have
%$$f_o(G)=\sum_{N\in {\cal P}}f_o(N)\ge \sum_{N\in {\cal P}} n_N/\chi(N)\ge n/\chi(G).$$
%\end{proof}

Now we will show that every connected chordal graph $G$ with even number of vertices satisfies Corollary \ref{cor:Scott}. As chordal graphs are perfect and their tree-width equals the order of maximum clique minus 1, we have $\chi(G)=\omega(G)=\tw(G)+1.$
For every graph $L$, $\chi(L)\le \tw(L)+1$ \cite{Dies}. Hence, for a minor $H$ of $G$, we have $\chi(H)\le \tw(H)+1\le \tw(G)+1=\chi(G)$. Hence $\mchi(G)=\chi(G)$ and so $G$ satisfies Conjecture \ref{conj:scott}. To construct a non-claw-free chordal graph
start from a complete graph $K_p$, choose $x\in V(K_p)$ and add four new vertices adjacent to $x$.
%Note that the main ingredient of the above proof is $\chi(G)=\tw(G)+1.$ Thus, together with chordal graphs without isolated vertices, odd cycles and many other graph classes whose graphs are obtained by clique-sum constructions, satisfy Conjecture \ref{conj:scott}. 






%\begin{lemma}\label{lem:RB}
%Let $F$ be a forest, and let $\ell:\ V(F)\to\{0,1\}$. Suppose that  $\sum_{v\in V(C)} \ell(v)$ is even for every connected component $C$ of $F$.
%Then there exists a bipartition
%\( V(F)=R\cup B\)
%such that both $F[R]$ and $F[B]$ are $\ell$-admissible in $F$.
%\end{lemma}
%\begin{proof}
%Without loss of generality, we may assume that $F$ is a tree. By Theorem \ref{thm:GY}, $V(F)$ can be partitioned into sets $V_1,\dots ,V_p$ such that each $F[V_i]$ is $\ell$-admissible. If we contract each $V_i$ into a vertex, we obtain a tree whose vertices can be properly colored into two colors. Thus, in $F$ we can color all vertices in each $V_i$ into one of two colors so that if the colors in $V_i$ and $V_j$, $1\le i\ne j\le p$, coincide then there is no edge between $V_i$ and $V_j$ in $F$. This implies 
%that if $R$ and $B$ are both the unions of vertices of $F$ of the same color, then they satisfy the conditions of the lemma. 
%\end{proof}
%\begin{theorem}
%Let $T$ be a tree of order $n$. Then
%\(
%f_{oe}(T)\ge \left\lfloor \frac n2\right\rfloor . 
%\)
%This bound is tight.
%Equivalently, for every labeling $\ell:V(T)\to\{0,1\}$, we have
%\(
%h_\ell(T)\ge \left\lfloor \frac n2\right\rfloor .
%\)
%\end{theorem}



\section{Conjecture on lower bound for $f_{oe}(G)/f_o(G)$}\label{sec:conj}

We state the following conjecture and provide two results supporting it. 

\begin{conjecture}\label{conj:half}
For every graph $G$, we have $f_{oe}(G)\ge f_o(G)/2$.
\end{conjecture}


%Note that the all-odd labeling $\ell\equiv 1$ is a special case, so the optimal universal constant in the labeled problem cannot exceed that in the all-odd problem. The following proposition makes this precise.

%\begin{proposition}\label{prop:constants-order}
%Let $c^*:=\sup\{c: f_o(H)\ge c|V(H)|\text{ for every graph }H\text{ without isolated vertices}\}$ and $\alpha^*:=\sup\{\alpha: h_{\ell}(G)\ge\alpha|V(G)|\text{ for all }G \text{      without isolated vertices},\text{ and all }\ell\}$. Then $\alpha^*\le c^*$.
%\end{proposition}
%\begin{proof}
%The labeling $\ell\equiv 1$ reduces the labeled problem to the all-odd problem.
%\end{proof}

\begin{proposition}\label{prop:tree}
For every tree $T$, we have $f_{oe}(T)\ge f_o(T)/2.$
\end{proposition}
\begin{proof}
Let $T$ have $n$ vertices and $\ell$ be an arbitrary labeling of $V(T)$.
By (\ref{ineq:1}),
$h_{\ell}(T) \ge \lceil{\frac{n-1}{2}}\rceil = \lfloor{\frac{n}{2}}\rfloor$.  If $T$ itself is odd, then $n$ is even and
$f_o(T)=n$, so
$h_{\ell}(T) \ge \frac{n}{2} = \frac{f_o(T)}{2}.$
If $T$ is not odd, then $f_o(T)\le n-1$. Note that $h_{\ell}(T) \ge \lfloor{\frac{n}{2}}\rfloor\ge \frac{n-1}{2}.$ 
Hence, $h_{\ell}(T) \ge \frac{f_o(T)}{2}.$
Since $\ell$ is arbitrary, we conclude that $f_{oe}(T)\ge f_o(T)/2$.
\end{proof}

%\begin{proposition}\label{prop:clique}
%For the complete graph $K_n$ with $n\ge 2$,$
%$$f_{oe}(K_n)=\Bigl\lfloor\frac{n}{2}\Bigr\rfloor.$$
%Moreover, $f_o(K_n)=n$ if $n$ is even and $f_o(K_n)=n-1$ if $n$ is odd, so in particular $f_{oe}(K_{2m})=\frac{1}{2}f_o(K_{2m})$.
%\end{proposition}
%\begin{proof} Fix a labeling $\ell$, and let $a=|L\cap \ell^{-1}(0)|$, $b=|L\cap \ell^{-1}(1)|$, so $a+b=n$.
%Every induced subgraph of $K_n$ on $t$ vertices is a clique $K_t$ in which every vertex has degree $t-1$. For an $\ell$-admissible $K_t$, either $t$ is odd (so $t-1$ is even) and all $t$ chosen vertices are labeled 0, or $t$ is even (so $t-1$ is odd) and all $t$ chosen vertices are labeled 1. Hence $h_{\ell}(K_n)=\max\{o(a),e(b)\}$, where $o(a)$ is the largest odd integer $\le a$ and $e(b)$ the largest even integer $\le b$.

%\emph{Lower bound.} Write $n=2m$ or $n=2m+1$. In both cases, it can be checked that for any split $a+b=n$, at least one of $o(a),e(b)$ is at least $m=\lfloor n/2\rfloor$.

%\emph{Upper bound.} For $n=2m$: take $a=b=m$; then $\max\{o(m),e(m)\}=m$. For $n=2m+1$: choose $a,b$ of appropriate parities to achieve $\max\{o(a),e(b)\}=m$.

%Hence $f_{oe}(K_n)=\lfloor n/2\rfloor$.
%\end{proof}

\begin{theorem}
 If $G$ is a complete multipartite graph, then $f_{oe}(G) \geq f_o(G)/2$.
\end{theorem}
\begin{proof}
%Let $n$ denote the order of $G$. 
Let $U \subseteq V(G)$ such that all degrees in $G[U]$ are odd and $|U|=f_o(G)$.  %Let $G_U = G[U]$ and 
Let $V_1,V_2,\ldots,V_p$ be the partite sets of $G[U]$. 
%and let $n_U=|U|$.
Note that $|U|$ is even as every vertex degree in $G[U]$ is odd. Since the degree in $G[U]$ of every vertex in $V_i$ is $|U|-|V_i|,$ which is odd, $|V_i|$ is odd for every $i \in [p]$. 
Since $|U|$ is even and each $|V_i|$ is odd, $p$ is even. 

Let $\ell : V(G) \rightarrow \{0,1\}$ be a labeling of $V(G)$ and let $V_i^j=\ell^{-1}(j)\cap V_i,$ where $j\in \{0,1\}.$
%\{ w \; | \; w \in V_i \; \& \; \phi(w)=j\}$. 
Define $r_1,r_2,\ldots,r_p$ such that $|V_i^{r_i}| > |V_i^{1-r_i}|$ for every $i \in [p]$, which is possible as $|V_i|$ is odd.
Let $W = \cup_{i=1}^p V_i^{r_i}$ and let $V_i^W = V_i \cap W$ for all $i \in [p]$.
As $|V_i^{r_i}| \geq |V_i^{1-r_i}|+1$, we have $|W| \geq (|U| + p)/2$.

Note that all vertices in $V_i^W$ have the same degree and also the same $\ell$-value. We call $V_i^W$ {\em good} if the degree and $\ell$-value are of the same parity for every vertex in $V_i^W$. Otherwise, we call $V_i^W$ {\em bad}.
Without loss of generality, assume that all $V_1^W,V_2^W,\ldots , V_s^W$ are bad and all $V_{s+1}^W,V_{s+2}^W,\ldots , V_p^W$ are good. 
We consider the following cases, where in Cases 1-3, $s$ is even and in Cases 4-6, $s$ is odd. 

\2

{\bf Case 1:} $s$ is even and $s \leq p/2$. We remove one vertex from each set $V_1^W,V_2^W,\ldots , V_s^W$ in $W$ and denote the resulting set $W^*$. Observe that $|W^*| \geq |W|- p/2 \geq |U| / 2$. Also, $W^*$ is $\ell$-admissible. Indeed, if $x\in W^*$ is in a good set, $d_{G[W^*]}(x)$ and $d_W(x)$ are of the same parity  since $s$ is even, implying that $d_{G[W^*]}(x)$ and $\ell(x)$ are of the same parity, and if $x\in W^*$ is in a bad set, $d_{G[W^*]}(x)$ and $d_{G[W]}(x)$ are of different parity since $s-1$ is odd, which means that $d_{W^*}(x)$ and $\ell(x)$ are of the same parity. 

\2

{\bf Case 2:} $s$ is even, $s > p/2$ and  $|V_i^{r_i}| \geq |V_i^{1-r_i}|+2$ for all $i \in [s]$. 
We remove one vertex from each set $V_1^W,V_2^W,\ldots , V_s^W$ in $W$ and denote the resulting set $W^*$.
We then have $|W^*| \geq |U| / 2$, as $W^*$ contains at least half of the vertices in $V_i$ for all $i \in [p]$. 
As in Case 1, $W^*$ is $\ell$-admissible.

\2

{\bf Case 3:} $s$ is even and $s > p/2$ and $|V_{i_0}^{r_{i_0}}| = |V_{i_0}^{1-r_{i_0}}|+1$ for some $i_0 \in [s]$. Since $p$ is even, $s \ge p/2+1$ and $p-s \le p/2-1$.
In $W$, we remove one vertex from each set $V_{s+1}^W,V_{s+2}^W,\ldots , V_p^W$ and replace  $V_{i_0}^{r_{i_0}}$ with $V_{i_0}^{1-r_{i_0}}$. We denote the resulting set $W^*$.
We have $|W^*| \geq |W| - (p/2-1)-1 \geq |U| / 2$. Note that for a vertex $x$ in $G[U]$ removing one vertex from another partite set flips parity the degree of $x$; 
removing one from the same partite set does not; replacing $V_{i_0}^{r_{i_0}}$ with $V_{i_0}^{1-r_{i_0}}$  changes the label and changes selected size by one in the margin-one case.
Hence, $W^*$ is $\ell$-admissible. 



\2

{\bf Case 4:} $s$ is odd and $s \geq p/2$. We remove one vertex from each set $V_{s+1}^W,V_{s+2}^W,\ldots , V_p^W$ in $W$ and denote the resulting set $W^*$. We note that again $W^*$ is $\ell$-admissible and $|W^*| \geq |W| - p/2 \geq |U|/ 2$.

\2

{\bf Case 5:} $s$ is odd and $s < p/2$ and  $|V_i^{r_i}| \geq |V_i^{1-r_i}|+2$ for all $i \in \{s+1,s+2,\ldots,p\}$. 
We remove one vertex from each set $V_{s+1}^W,V_{s+2}^W,\ldots , V_p^W$ in $W$ and denote the resulting set $W^*$.
We then have $|W^*| \geq |U| / 2$, as $W^*$ contains at least half of the vertices in $V_i$ for all $i \in [p]$. 
As before, $W^*$ is $\ell$-admissible.

\2

{\bf Case 6:} $s$ is odd, $s < p/2$ and $|V_{i_0}^{r_{i_0}}| = |V_{i_0}^{1-r_{i_0}}|+1$ for some ${i_0} \in \{s+1,s+2,\ldots,p\}$. Since $p$ is even, $s \le p/2-1$.
Then remove one vertex from each set $V_{1}^W,V_{2}^W,\ldots , V_s^W$ and replace $V_{i_0}^{r_{i_0}}$ with $V_{i_0}^{1-r_{i_0}}$. We denote the resulting set $W^*$.
Then $|W^*| \geq |W| - (p/2-1)-1 \geq |U| / 2$ and, similar to Case 3, $W^*$ is $\ell$-admissible. 

Since $\ell$ is arbitrary, we conclude that $f_{oe}(G) \geq f_o(G)/2$.
\end{proof}




%\paragraph{Acknowledgement} The authors used ChatGPT Pro during an early stage of exploring how to directly prove Corollary \ref{cor:RB}. 


\begin{thebibliography}{99}
\bibitem{Barioli+} F. Barioli, W. Barrett, S. Fallat, H. T. Hall, L. Hogben, B. Shader, P. van den Driessche, and H. van der Holst, Parameters related to tree-width, zero forcing, and maximum nullity of a graph, J. Graph Theory 72 (2013), 146--177.
\bibitem{Caro1994} Y. Caro, On induced subgraphs with odd degrees, Discrete Math. 132 (1994), 23--28.
%\bibitem{CKR} Y. Caro, I. Krasikov and Y. Roditty, On induced subgraphs of trees with restricted degrees, 
\bibitem{CLZ} Y. Caro, J. Lauri and C. Zarb. Two short proofs of the perfect forest
theorem, Theory and Applications of Graphs 4(1) (2017), article 4.
\bibitem{DP} M. Delcourt and L. Postle, Reducing linear Hadwiger's conjecture to coloring small graphs, J. Amer. Math. Soc. 38(2) (2025), 481--507.
\bibitem{Dies} R. Diestel, Graph Theory, 6th edition, Springer, 2025.
\bibitem{FK} A. Ferber and M. Krivelevich, Every graph contains a linearly sized induced subgraph with all degrees odd, Adv. Math. 406 (2022) 108534.
\bibitem{Gutin2016} G. Gutin, Note on perfect forests, J. Graph Theory 82 (2016), no. 3, 233--235.
\bibitem{GutinYeo2022} G. Gutin and A. Yeo, Perfect forests in graphs and their extensions, J. Graph Theory 101(3) (2022), 572--592. 
\bibitem{Halin} R. Halin, $S$-functions for graphs, J. Geometry 8 (1976), 171--186.
\bibitem{Hogben} L. Hogben, Nordhaus–Gaddum problems for Colin de Verdi{\'e}re type parameters, variants of tree-width, and related parameters, 
in Recent Trends in Combinatorics, Springer, 2016, pp. 275--294.
\bibitem{HYLL} X. Hou, L. Yu, J. Li, and B. Liu, Odd Induced Subgraphs in Graphs with Treewidth at Most Two, Graphs Combin. 34 (2018), 535--544.
%The IMA Volumes in Mathematics and its Applications, vol. 159, 
\bibitem{L1993} L. Lov\'asz, Combinatorial Problems and Exercises, 2nd edition, AMS Chelsea Publishing, 1993.
\bibitem{Ning} B. Ning, On Scott’s odd induced subgraph conjecture and a related problem, arXiv:2604.19727v1, 2026.
\bibitem{RS} A.J. Radcliffe and A.D. Scott, Every tree contains a large induced subgraph with all degrees odd, Discrete Math 140(1--3) (1995),  275--279.
%\bibitem{RHZ} M. Rao, J. Hou, and Q. Zeng, Odd induced subgraphs in planar graphs with large girth, Graphs Combin. 38 (2022), Paper No. 105.
%\bibitem{Schr} A. Schrijver, Theory of Linear and Integer Programming, John Wiley \& Sons, 1998.
\bibitem{Scott1992} A.D. Scott, Large induced subgraphs with all degrees odd, Comb. Probab. Comput. 1 (1992) 335--349.
\bibitem{Scott2001}  A.D. Scott, On induced subgraphs with all degrees odd, Graphs Combin.
17(3) (2001) 539--553.
\bibitem{WW} T. Wang and B. Wu, Maximum odd induced subgraph of a graph concerning its chromatic number, J. Graph Theory 107 (2024), 
578--596.
\end{thebibliography}

\end{document}

It is not hard to obtain the following corollaries, which show that the bound for $f_{oe}(G)$ in Theorem \ref{thm:tw} is sharp. 

\begin{corollary}\label{cor:planar}
If $G$ is a connected planar graph of order $n$, then $f_{oe}(G)\ge \lceil (n-1)/4 \rceil$ and this bound is sharp.
\end{corollary}
\begin{proof}
The bound $f_{oe}(G)\ge \lceil (n-1)/4 \rceil$ follows from the Four Color Theorem and the fact that planarity is closed under taking minors \cite{Dies}. The bound is sharp because
$f_{oe}(K_{2,2,2})\ge 2$ and $f_o(K_{2,2,2})=2.$ To show that $f_o(K_{2,2,2})=2$, consider a maximum order induced odd subgraph $H$ of $K_{2,2,2}$ and let $x_1,x_2,x_3$ be the number of vertices of $H$ in each partite set of $K_{2,2,2}$. If $\min\{x_1,x_2,x_3\}=0$ then 
we observe that $f_o(K_{2,2})=2.$ Otherwise, for each $i\in [3]$, we have $x_i>0$ and $|V(H)|-x_i$ is odd. Now, if $|V(H)|$ is odd,
then each $x_i=2$, a contradiction with $|V(H)|$ being odd. If $|V(H)|$ is even then each $x_i=1$, a contradiction with $|V(H)|$ being even.
\end{proof}

\begin{corollary}\label{cor:RB}
Let $T$ be an $n$-vertex tree. Then $f_{oe}(T)\ge \lfloor n/2 \rfloor$ and this bound is sharp. 
\end{corollary}
\begin{proof}
The bound of Corollary \ref{cor:RB} follows from Theorem \ref{thm:tw} and the fact that $\lceil \frac{n-1}{2}\rceil=\lfloor \frac{n}{2}\rfloor$. To show sharpness of this bound it enough to see that it becomes an equality already for $K_2$: label both vertices by 0. 
\end{proof}

In fact, it is not hard to prove that 
\(
f_{oe}(K_{1,m})=\left\lfloor \frac{m+1}{2}\right\rfloor 
\)
for each $m\ge 1.$

\begin{proposition}
For the star \(K_{1,m}\), where \(m\ge 1\),
\[
f_{oe}(K_{1,m})=\left\lfloor \frac{m+1}{2}\right\rfloor .
\]
\end{proposition}
\begin{proof}
Let \(c\) be the center of \(K_{1,m}\), and let \(L\) be the set of its \(m\) leaves.
Fix a labeling
$
\ell:V(K_{1,m})\to\{0,1\}
$ and 
let
$
a=|L\cap \ell^{-1}(0)|, b=|L\cap \ell^{-1}(1)|,
$
so that \(a+b=m\).

We first determine \(h_\ell(K_{1,m})\). Let \(U\subseteq V(K_{1,m})\) be
\(\ell\)-admissible.

If \(c\notin U\), then every leaf in $U$ has degree \(0\) in \(K_{1,m}[U]\).
Thus, $ |U|\le a.$
Conversely, the set of all \(0\)-labeled leaves is \(\ell\)-admissible, so this
case the maximum of  \(|U|\) is $a$.

Now suppose \(c\in U\). Since every leaf in $U$ has degree \(1\) in
\(K_{1,m}[U]\), every such leaf must be \(1\)-labeled. If \(t\) denotes the
number of such leaves, then the center has degree \(t\) in \(K_{1,m}[U]\),
so we also require
$
t\equiv \ell(c)\pmod 2.
$
Thus the maximum of $|U|$ is
\(
1+t_\ell,
\)
where \(t_\ell\) is the largest integer satisfying
\(
0\le t_\ell\le b, t_\ell\equiv \ell(c)\pmod 2,
\)
if such an integer exists. If no such integer exists, then no \(\ell\)-admissible
induced subgraph containing the center exists. The only exceptional case is
\(b=0\) and \(\ell(c)=1\). Hence
\[
h_\ell(K_{1,m})=\max\{a,\ 1+t_\ell\},
\]
where the second term equals \(0\) when \(b=0\) and \(\ell(c)=1\).

We now prove the lower bound $h_\ell(K_{1,m})\ge \left\lfloor\frac{m+1}{2}\right\rfloor.$ If
\(
a\ge \left\lfloor\frac{m+1}{2}\right\rfloor,
\)
then the set of all \(0\)-labeled leaves is already an \(\ell\)-admissible set
of size at least \(\left\lfloor\frac{m+1}{2}\right\rfloor\).

Otherwise,
\(
a<\left\lfloor\frac{m+1}{2}\right\rfloor.
\)
Since \(a+b=m\), this implies
\(
b\ge \left\lceil\frac{m+1}{2}\right\rceil .
\)
Among the two largest integers \(b\) and \(b-1\), one has parity \(\ell(c)\).
Therefore there exists an integer \(t\le b\) with
\(
t\equiv \ell(c)\pmod 2
\)
and
\(
t\ge b-1.
\)
Thus the maximum size of $U$ containing \(c\) is at least
\[
1+(b-1)=b\ge \left\lceil\frac{m+1}{2}\right\rceil
\ge \left\lfloor\frac{m+1}{2}\right\rfloor .
\]
Therefore, for every labeling \(\ell\),
\[
h_\ell(K_{1,m})\ge \left\lfloor\frac{m+1}{2}\right\rfloor .
\]
Hence
\(
f_{oe}(K_{1,m})\ge \left\lfloor\frac{m+1}{2}\right\rfloor .
\)

It remains to prove the matching upper bound. We exhibit labelings for which no larger
\(\ell\)-admissible induced subgraph exists.
First suppose \(m=2r\). Label \(r\) leaves by \(0\) and \(r\) leaves by
\(1\). Choose the label of the center so that
\(
\ell(c)\not\equiv r\pmod 2.
\)
Then an \(\ell\)-admissible set not containing the center has size at most \(r\), while
an \(\ell\)-admissible set containing the center can contain at most \(r-1\) of the
\(1\)-labeled leaves, and therefore has size at most \(1+(r-1)=r\). Hence
\[
h_\ell(K_{1,2r})\le r=\left\lfloor\frac{2r+1}{2}\right\rfloor .
\]

Now suppose \(m=2r+1\). Label \(r+1\) leaves by \(0\) and \(r\) leaves
by \(1\). Choose the label of the center so that
$
\ell(c)\not\equiv r\pmod 2.
$
Then an \(\ell\)-admissible set not containing the center has size at most \(r+1\),
while an \(\ell\)-admissible set containing the center can contain at most \(r-1\)
\(1\)-labeled leaves, and therefore has size at most \(r\). Hence
\[
h_\ell(K_{1,2r+1})\le r+1
=
\left\lfloor\frac{2r+2}{2}\right\rfloor .
\]
Combining the lower and upper bounds gives
\[
f_{oe}(K_{1,m})=\left\lfloor \frac{m+1}{2}\right\rfloor .
\]
\end{proof}


\section{Generalizing Gallai's Theorem}\label{sec:Gallai}

In this section we characterize digraphs which have even-even or even-odd partitions, and unify both problems into a single theorem with arbitrary side-dependent parity prescriptions. We use linear algebra over $\Ftwo$, leading to algorithms with running time $O(|V|^3)$ arithmetic operations.

Let $M$ be the adjacency matrix of $D=(V,A)$ over $\Ftwo$, i.e.\ $M_{uv}=1$ if $uv\in A$ and 0 otherwise. Let $p=(p_v)_{v\in V}$ be the parity vector of $D$, where $p_v=1$ if $d^+(v)$ is odd and $p_v=0$ otherwise.

\subsection{Unified parity-prescription theorem}

For vectors $a,b\in\Ftwo^V$, we seek a partition $V=V_0\cup V_1$ such that
$$
\deg^+_{D[V_0]}(v)\equiv a_v\pmod{2} \text{ for } v\in V_0,
\quad
\deg^+_{D[V_1]}(v)\equiv b_v\pmod{2} \text{ for } v\in V_1.
$$

\begin{theorem}\label{thm:general-digraph}
Let $D=(V,A)$ be a digraph with adjacency matrix $M$ over $\Ftwo$ and parity vector $p$. Let $a,b\in\Ftwo^V$. Then there exists a partition $V=V_0\cup V_1$ with the prescribed parities if and only if the linear system
\begin{equation}\label{eq:general-system}
\bigl(M+\diaf_{oe}(p+a+b)\bigr)s = p+a
\end{equation}
is solvable over $\Ftwo$, where $s\in\Ftwo^V$ is the indicator vector of $V_1$. Every solution $s$ yields such a partition.
\end{theorem}
\begin{proof}
Fix $s\in\Ftwo^V$ and let $P_v(s)$ denote the parity of the number of out-neighbors of $v$ lying in the same part as $v$. Using the identity $[s_u=s_v]=1+s_u+s_v$ over $\Ftwo$,
$$
P_v(s)\equiv \sum_{u:vu\in A}(1+s_u+s_v)\equiv p_v+(Ms)_v+p_vs_v\pmod{2}.
$$
The desired parity at $v$ is $a_v$ when $s_v=0$ and $b_v$ when $s_v=1$, which can be written uniformly as $a_v+(a_v+b_v)s_v$ over $\Ftwo$. Setting $P_v(s)$ equal to this target and rearranging gives
$$
(Ms)_v+(p_v+a_v+b_v)s_v=p_v+a_v,
$$
which is the $v$-th coordinate of \eqref{eq:general-system}. Hence $s$ defines a valid partition if and only if it satisfies \eqref{eq:general-system}.
\end{proof}

The even-even and even-odd problems are immediate special cases.

\begin{corollary}\label{cor:original-digraph}
Let $D=(V,A)$ be a digraph with adjacency matrix $M$ and parity vector $p$.
\begin{enumerate}
\item $D$ has an even-even partition if and only if $(M+\diaf_{oe}(p))s=p$ is solvable over $\Ftwo$.
\item $D$ has an even-odd partition if and only if $(M+I+\diaf_{oe}(p))s=p$ is solvable over $\Ftwo$.
\end{enumerate}
In both cases, deciding whether such a solution exists and finding one can be done in $O(|V|^3)$ arithmetic operations over $\Ftwo$.
\end{corollary}
\begin{proof}
For even-even take $a=b=0$; for even-odd take $a=0$ and $b=\mathbf{1}$. The running time follows from Gaussian elimination over $\Ftwo$.
\end{proof}

The general system also yields an exact count of all valid partitions.

\begin{corollary}\label{cor:counting}
Fix $a,b\in\Ftwo^V$ and let $A_0=M+\diaf_{oe}(p+a+b)$. If the system $A_0 s=p+a$ is inconsistent, no valid partition exists. If it is consistent, the number of valid partitions is
$$
2^{|V|-\operatorname{rank}(A_0)}.
$$
\end{corollary}
\begin{proof}
The solution set of a consistent linear system over a field is an affine subspace of dimension $|V|-\operatorname{rank}(A_0)$.
\end{proof}

\subsection{Structural results}

We make the following observation on a large family of digraphs that always admit both types of partition.

\begin{remark}\label{rem:construction}
Let $S=(W,B)$ be a symmetric digraph (i.e.\ $xy\in B$ iff $yx\in B$). By Gallai's Theorem, $S$ has an even-even partition $W=W_0\cup W_1$ and an even-odd partition $W=W'_0\cup W'_1$. Let $L=(U,C)$ be a digraph with $U\cap W=\varnothing$ in which every vertex has even out-degree. Obtain a new digraph $D$ by adding arcs from $U$ to $W$ such that every $x\in U$ sends an even number of arcs to both $W_0$ and $W'_0$. Then $D$ has an even-even partition with parts $W_0\cup U$ and $W_1$, and an even-odd partition with parts $W'_0\cup U$ and $W'_1$.
\end{remark}

The following proposition generalizes this remark into a reusable extension lemma.

\begin{proposition}\label{prop:extension}
Let $S=(W,B)$ and $L=(U,C)$ be vertex-disjoint digraphs, and let $D$ be obtained by adding arcs from $U$ to $W$ only.
\begin{enumerate}
    \item Suppose $W=X\cup Y$ is an even-even partition of $S$, every vertex of $L$ has even out-degree in $L$, and every vertex of $U$ sends an even number of arcs to $X$. Then $(X\cup U, Y)$ is an even-even partition of $D$.
    \item Suppose $W=X\cup Y$ is an even-odd partition of $S$ with $X$ the even side, every vertex of $L$ has even out-degree in $L$, and every vertex of $U$ sends an even number of arcs to $X$. Then $(X\cup U, Y)$ is an even-odd partition of $D$.
\end{enumerate}
\end{proposition}
\begin{proof}
For $u\in U$, its out-neighbors in the part $X\cup U$ are its out-neighbors in $U$ together with those in $X$; both counts are even by assumption. For $w\in W$, no arcs go from $W$ to $U$, so the out-neighbors of $w$ inside its own part are the same in $D$ as in $S$, and every parity condition already satisfied in $S$ remains satisfied in $D$.
\end{proof}

\begin{proposition}\label{prop:bipartite-symmetric}
Let $D$ be a digraph.
\begin{enumerate}
    \item If the underlying undirected graph of $D$ is bipartite, then $D$ has an even-even partition.
    \item If $D$ is symmetric, then $D$ has both an even-even partition and an even-odd partition.
\end{enumerate}
\end{proposition}
\begin{proof}
(1) Let $V=X\cup Y$ be a bipartition. Every arc goes between $X$ and $Y$, so $D[X]$ and $D[Y]$ are edgeless, giving out-degrees 0 (even) everywhere.

(2) Since $D$ is symmetric, the out-degree of $v$ in $D[X]$ equals the degree of $v$ in the underlying undirected graph $G[X]$ for any $X$. Any partition provided by Gallai's Theorem for $G$ therefore yields the corresponding even-even or even-odd partition for $D$.
\end{proof}

We conclude with a complete characterization for functional digraphs (out-degree at most 1).

\begin{proposition}\label{prop:functional}
Let $D=(V,A)$ be a digraph in which every vertex has out-degree at most $1$.
\begin{enumerate}
    \item $D$ has an even-even partition if and only if every directed cycle of $D$ has even length.
    \item $D$ has an even-odd partition if and only if every vertex of out-degree $0$ is isolated.
\end{enumerate}
\end{proposition}
\begin{proof}
\textbf{(1).} Suppose $D$ has an even-even partition $V=V_0\cup V_1$. For any arc $uv$, if $u$ and $v$ were in the same part then $u$ would have odd out-degree 1 inside that part, contradicting the even-even condition. Hence every arc crosses the partition, so the underlying undirected graph is bipartite and all cycles have even length. In a digraph with out-degrees at most 1, every undirected cycle is directed cyclically, so every directed cycle has even length.

Conversely, if every directed cycle has even length then the underlying undirected graph is bipartite (since its cycles are exactly the directed cycles). A bipartition gives an even-even partition by part (1) of Proposition~\ref{prop:bipartite-symmetric}.

\textbf{(2).} Suppose $D$ has an even-odd partition $V=V_0\cup V_1$. Let $v$ have out-degree 0. Then $v$ has out-degree 0 in every induced subgraph, so $v$ cannot be in $V_1$ (where odd out-degree is required), hence $v\in V_0$. Now let $u\to v$ be an arc with $v\in V_0$; then $u$'s only out-neighbor lies in $V_0$, giving out-degree 0 inside $V_1$ if $u\in V_1$ and out-degree 1 (odd) inside $V_0$ if $u\in V_0$. Both cases violate the partition condition. Hence $v$ has no in-neighbors either, so $v$ is isolated.

Conversely, suppose every vertex of out-degree 0 is isolated. Let $V_0$ be the set of isolated vertices and $V_1:=V\setminus V_0$. Every vertex in $V_1$ has out-degree exactly 1, and its unique out-neighbor must lie in $V_1$ (since vertices of $V_0$ are isolated). Therefore every vertex of $V_1$ has out-degree 1 inside $D[V_1]$, and every vertex of $V_0$ has out-degree 0 inside $D[V_0]$. This is an even-odd partition.
\end{proof}

